% ============================================================
% supplement.tex — Online Supplement
% Biomarker Dilution Simulation Study
% ============================================================
\documentclass[12pt]{article}

% --- Page layout ---
\usepackage[margin=1in]{geometry}
\usepackage{setspace}
\doublespacing

% --- Line numbering ---
\usepackage{lineno}
\linenumbers

% --- Mathematics ---
\usepackage{amsmath,amssymb}

% --- Graphics ---
\usepackage{graphicx}
\graphicspath{{figures/}}

% --- Tables ---
\usepackage{booktabs}
\usepackage{multirow}
\usepackage{array}
\usepackage{longtable}

% --- Colors ---
\usepackage[dvipsnames]{xcolor}

% --- Hyperlinks ---
\usepackage[colorlinks=true,linkcolor=blue,citecolor=blue,urlcolor=blue]{hyperref}

% --- Bibliography ---
\usepackage[numbers,sort&compress]{natbib}

% --- Fonts ---
\usepackage[T1]{fontenc}
\usepackage{mathptmx}

% --- Misc ---
\usepackage{textcomp}
\usepackage{url}

% --- Custom commands ---
\newcommand{\Xobs}{\mathbf{X}_{\mathrm{obs}}}
\newcommand{\Xtrue}{\mathbf{X}_{\mathrm{true}}}

% --- Supplement figure/table numbering ---
\renewcommand{\thefigure}{E\arabic{figure}}
\renewcommand{\thetable}{E\arabic{table}}

% ============================================================
\begin{document}

\begin{center}
{\Large\bfseries Online Supplement}

\vspace{0.5cm}

{\large Sample Dilution Systematically Distorts Biomarker Discovery in Lower Respiratory Tract Fluid: A Simulation Framework with Empirical Validation}

\vspace{0.5cm}

Aartik Sarma, MD

\end{center}

\vspace{1cm}

\tableofcontents

\newpage

% ============================================================
\section{Extended Methods}
% ============================================================

\subsection{Mathematical Details of Normalization Methods}

\subsubsection{Probabilistic Quotient Normalization (PQN)}

For a dataset $\mathbf{X} \in \mathbb{R}^{n \times p}$:
\begin{enumerate}
\item Compute the reference spectrum $\mathbf{r} \in \mathbb{R}^{p}$: $r_j = \mathrm{median}_i(x_{ij})$
\item For each sample $i$, compute quotients: $q_{ij} = x_{ij} / r_j$
\item Compute the normalization factor: $f_i = \mathrm{median}_j(q_{ij})$
\item Normalize: $x_{ij}^{\mathrm{norm}} = x_{ij} / f_i$
\end{enumerate}

\subsubsection{Centered Log-Ratio (CLR) Transformation}

For sample $i$ with biomarker values $\mathbf{x}_i = (x_{i1}, \ldots, x_{ip})$:
\begin{equation}
\mathrm{clr}(\mathbf{x}_i)_j = \ln(x_{ij}) - \frac{1}{p} \sum_{k=1}^{p} \ln(x_{ik}) = \ln\left(\frac{x_{ij}}{g(\mathbf{x}_i)}\right)
\end{equation}
where $g(\mathbf{x}_i) = \left(\prod_{k=1}^{p} x_{ik}\right)^{1/p}$ is the geometric mean.

\textbf{Key property:} The CLR transformation satisfies $\sum_{j=1}^{p} \mathrm{clr}(\mathbf{x}_i)_j = 0$ for all $i$. This zero-sum constraint is the source of the induced negative correlations discussed in the main text.

\textbf{Dilution invariance~\cite{Aitchison1986}:} If $\mathbf{x}_i^{\mathrm{obs}} = d_i \cdot \mathbf{x}_i^{\mathrm{true}}$, then:
\begin{equation}
\mathrm{clr}(\mathbf{x}_i^{\mathrm{obs}})_j = \ln\left(\frac{d_i x_{ij}^{\mathrm{true}}}{d_i \cdot g(\mathbf{x}_i^{\mathrm{true}})}\right) = \ln\left(\frac{x_{ij}^{\mathrm{true}}}{g(\mathbf{x}_i^{\mathrm{true}})}\right) = \mathrm{clr}(\mathbf{x}_i^{\mathrm{true}})_j
\end{equation}
demonstrating exact cancellation of the dilution factor under proportional dilution.

\subsubsection{Quantile Normalization}

\begin{enumerate}
\item Sort each sample's values independently
\item Replace sorted values with the average of the sorted values at each rank across all samples
\item Re-order to original positions
\end{enumerate}

This forces identical marginal distributions across all samples.

\subsection{RV Coefficient}

The RV coefficient~\cite{Robert1976} between two data matrices $\mathbf{X}_1$ and $\mathbf{X}_2$ is defined as:
\begin{equation}
\mathrm{RV}(\mathbf{X}_1, \mathbf{X}_2) = \frac{\mathrm{tr}(\mathbf{S}_{12} \mathbf{S}_{21})}{\sqrt{\mathrm{tr}(\mathbf{S}_{11}^2) \cdot \mathrm{tr}(\mathbf{S}_{22}^2)}}
\end{equation}
where $\mathbf{S}_{kl} = \mathbf{X}_k^T \mathbf{X}_l / n$ and $\mathrm{tr}$ denotes the matrix trace. The RV coefficient ranges from 0 (no association) to 1 (identical subspaces).

\subsection{Adjusted Rand Index}

The adjusted Rand index (ARI)~\cite{Hubert1985} is defined as:
\begin{equation}
\mathrm{ARI} = \frac{\mathrm{RI} - \mathbb{E}[\mathrm{RI}]}{\max(\mathrm{RI}) - \mathbb{E}[\mathrm{RI}]}
\end{equation}
where RI is the Rand index (fraction of pairs of samples that are either in the same cluster in both partitions or in different clusters in both partitions). The ARI corrects for chance agreement and ranges from $-1$ to 1, with 0 indicating random clustering.

\subsection{Jaccard Similarity}

For two feature sets $A$ and $B$:
\begin{equation}
J(A, B) = \frac{|A \cap B|}{|A \cup B|}
\end{equation}

\subsection{Simulation Parameter Justification}

\subsubsection{Beta Distribution for Dilution}

The Beta distribution was selected because: (1) it is bounded on $[0, 1]$, reflecting that dilution factors are proportions; (2) it is highly flexible, accommodating symmetric, left-skewed, and right-skewed distributions through two shape parameters; (3) the parameterization allows intuitive interpretation, with the mean $\alpha/(\alpha+\beta)$ directly representing the expected proportion of signal retained.

The specific parameterizations were chosen to span the range of clinically observed BAL recovery rates: 60--95\% fluid recovery (mild), 30--70\% (moderate), and 5--40\% (severe).

\subsubsection{Log-Normal Distribution for Biomarkers}

Biomarker concentrations in biological fluids are typically right-skewed with occasional extreme values. The log-normal distribution naturally produces these characteristics and is the standard generative model in metabolomics and proteomics simulation studies.

\subsubsection{Effect Size Selection}

Cohen's $d = 0.8$ was chosen as the primary effect size because it represents a ``large'' effect in behavioral science terminology but corresponds to realistic fold changes ($\sim$1.5--2.0) commonly targeted in BAL biomarker studies.

% ============================================================
\section{Extended Results}
% ============================================================

\subsection{High-Dimensional Analysis (Figure~E1)}

\begin{table}[ht!]
\centering
\caption{Performance Metrics in High-Dimensional Settings Under Moderate Dilution with CLR Normalization}
\label{tab:highdim}
\begin{tabular}{lccccc}
\toprule
$p$ & Power (raw) & Power (FDR) & Type~I (raw) & Type~I (FDR) & Comp.\ Time (s) \\
\midrule
10 & 0.994 & 0.980 & 0.783 & 0.320 & 2.1 \\
50 & 0.992 & 0.945 & 0.775 & 0.215 & 5.8 \\
100 & 0.990 & 0.910 & 0.770 & 0.180 & 12.4 \\
500 & 0.988 & 0.855 & 0.765 & 0.142 & 89.7 \\
\bottomrule
\end{tabular}
\end{table}

\subsection{Type~I Error Under Null (Figure~E2)}

Under complete null conditions (zero effect size for all $p = 10$ biomarkers), all normalization methods maintained nominal Type~I error rates regardless of test type (Table~\ref{tab:type1_null}). This demonstrates that CLR's Type~I error inflation (78--80\% in the primary simulation, Table~2 in main text) arises specifically when true differential biomarkers exist: the zero-sum constraint propagates real biological signals to null features, creating apparent differences where none exist biologically. Under pure null conditions, no such signal propagation occurs, and all methods behave nominally.

\begin{table}[ht!]
\centering
\caption{Type~I Error Under Complete Null Conditions (Zero Effect Size)}
\label{tab:type1_null}
\begin{tabular}{lcc}
\toprule
Normalization & $t$-Test Type~I Error & Permutation Type~I Error \\
\midrule
None & 0.055 & 0.050 \\
CLR & 0.035 & 0.035 \\
PQN & 0.030 & 0.030 \\
Total sum & 0.020 & 0.020 \\
Median & 0.030 & 0.040 \\
Quantile & 0.040 & 0.055 \\
\bottomrule
\end{tabular}
\end{table}

This result refines our mechanistic understanding: the compositional constraint inherent in CLR acts as a ``signal amplifier'' that redistributes variance from differential to non-differential features. When true signal exists, this redistribution creates false positives among null features. Permutation-based testing, by preserving the joint dependency structure under the null hypothesis, correctly accounts for this constraint and controls Type~I error at nominal levels even when true effects are present (see primary simulation results).

\subsection{Batch--Dilution Interaction Results (Figure~E3)}

\begin{table}[ht!]
\centering
\caption{Classification Accuracy Under Combined Batch Effects and Dilution}
\label{tab:batch_dilution}
\begin{tabular}{lcccc}
\toprule
Correction Strategy & Power & Type~I Err. & Clust.\ ARI & Class.\ Acc. \\
\midrule
None & 0.820 & 0.055 & 0.008 & 0.885 \\
Dilution only (CLR) & 0.985 & 0.780 & 0.890 & 0.975 \\
Batch only (ComBat) & 0.870 & 0.048 & 0.025 & 0.920 \\
Combined (ComBat + CLR) & 0.990 & 0.750 & 0.950 & 0.990 \\
Combined + permutation & 0.960 & 0.055 & 0.950 & 0.990 \\
\bottomrule
\end{tabular}
\end{table}

\subsection{PICFLU Reanalysis Details (Figure~E4)}

The PICFLU cohort BAL data from Boyd et al.~\cite{Boyd2020} comprised 84 critically ill children with influenza, with 49 cytokine/protein measurements in lower respiratory tract fluid. Clinical groups were defined by prolonged acute hypoxemic respiratory failure (PrAHRF) or death (22 cases vs.\ 62 controls). Total protein exhibited a coefficient of variation of 115.5\% (mean 25,438, SD 29,390 pg/mL), reflecting substantial specimen dilution variability. Fifteen of 49 analytes (31\%) correlated significantly with total protein ($p < 0.05$), with mean $|r| = 0.19$.

After applying each normalization method independently:

\begin{itemize}
\item \textbf{None (raw):} 13 significant analytes (nominal $p < 0.05$), 0 after FDR correction
\item \textbf{CLR:} 17 significant (nominal), 12 after FDR---the most significant analytes of any method
\item \textbf{Total sum:} 15 significant (nominal), 3 after FDR
\item \textbf{PQN:} 12 significant (nominal), 0 after FDR
\item \textbf{Median:} 12 significant (nominal), 0 after FDR
\item \textbf{Quantile:} 12 significant (nominal), 0 after FDR
\item \textbf{Concordance:} Average pairwise Jaccard similarity = 0.49 (range: 0.25--0.71)
\end{itemize}

These patterns directly mirror the simulation findings: CLR identified the most significant biomarkers (consistent with its high sensitivity), while unnormalized and median analyses were more conservative. The moderate inter-method Jaccard concordance confirms that normalization choice materially affects biomarker discovery in real BAL data.

\subsection{Multi-Group Extension}

When the number of groups was increased from 2 to 3 and 4 (using ANOVA/Kruskal--Wallis instead of $t$-tests), the qualitative patterns of normalization performance were preserved. Power decreased slightly with more groups (as expected with fixed total sample size), and clustering ARI differences between normalization methods were more pronounced with 3--4 groups.

\subsection{Limit of Detection Sensitivity Analysis (Figures~E6--E7)}

\subsubsection{Methods}

Many analytes in BAL and other biofluids yield measurements below the assay's lower limit of detection (LOD), particularly after dilution attenuates already-low concentrations. The choice of how to handle these censored values is rarely justified but can influence downstream results. We evaluated five imputation strategies across five LOD severity levels.

LOD severity was controlled by setting the LOD threshold at the specified percentile of each biomarker's diluted concentration distribution (5\%, 10\%, 20\%, 30\%, or 50\%). Five imputation methods were compared:
\begin{enumerate}
\item \textbf{Zero substitution}: Below-LOD values set to 0
\item \textbf{LOD/2}: Below-LOD values set to half the LOD
\item \textbf{LOD/$\sqrt{2}$}: Hornung and Reed substitution~\cite{Hornung1990}, the most common default
\item \textbf{LOD}: Below-LOD values set to the LOD value itself (upper bound)
\item \textbf{MICE}: Below-LOD values set to missing, then imputed via Multiple Imputation by Chained Equations~\cite{vanBuuren2011} using scikit-learn's \texttt{IterativeImputer} (max\_iter=10, min\_value=0)
\end{enumerate}

All 25 severity $\times$ imputation combinations were evaluated under moderate dilution (Beta(5,5)) across all 7 normalization methods with 20 replications each, totaling 3,500 simulation runs.

\subsubsection{Results}

Table~\ref{tab:lod_sensitivity} summarizes performance metrics averaged across normalization methods.

\begin{table}[ht!]
\centering
\caption{LOD Sensitivity: Performance Metrics by Severity and Imputation Method (Averaged Across Normalization Methods)}
\label{tab:lod_sensitivity}
\small
\begin{tabular}{llcccccc}
\toprule
LOD Sev. & Imputation & Power & Type~I & Corr. & PCA RV & ARI & Acc. \\
\midrule
\multirow{5}{*}{5\%}
 & Zero & 0.744 & 0.258 & 0.633 & 0.317 & 0.249 & 0.924 \\
 & LOD/2 & 0.903 & 0.464 & 0.780 & 0.417 & 0.495 & 0.983 \\
 & LOD/$\sqrt{2}$ & 0.900 & 0.460 & 0.785 & 0.428 & 0.505 & 0.983 \\
 & LOD & 0.899 & 0.465 & 0.789 & 0.436 & 0.515 & 0.982 \\
 & MICE & 0.897 & 0.460 & 0.789 & 0.429 & 0.413 & 0.944 \\
\midrule
\multirow{5}{*}{10\%}
 & Zero & 0.675 & 0.209 & 0.568 & 0.287 & 0.117 & 0.895 \\
 & LOD/2 & 0.896 & 0.450 & 0.779 & 0.413 & 0.476 & 0.982 \\
 & LOD/$\sqrt{2}$ & 0.898 & 0.464 & 0.786 & 0.430 & 0.508 & 0.982 \\
 & LOD & 0.899 & 0.469 & 0.791 & 0.442 & 0.509 & 0.981 \\
 & MICE & 0.885 & 0.431 & 0.791 & 0.419 & 0.320 & 0.892 \\
\midrule
\multirow{5}{*}{20\%}
 & Zero & 0.698 & 0.217 & 0.562 & 0.297 & 0.100 & 0.878 \\
 & LOD/2 & 0.897 & 0.458 & 0.784 & 0.416 & 0.478 & 0.981 \\
 & LOD/$\sqrt{2}$ & 0.898 & 0.469 & 0.792 & 0.438 & 0.496 & 0.980 \\
 & LOD & 0.899 & 0.462 & 0.796 & 0.454 & 0.495 & 0.975 \\
 & MICE & 0.848 & 0.360 & 0.779 & 0.385 & 0.168 & 0.778 \\
\midrule
\multirow{5}{*}{30\%}
 & Zero & 0.737 & 0.233 & 0.576 & 0.308 & 0.095 & 0.874 \\
 & LOD/2 & 0.896 & 0.452 & 0.778 & 0.419 & 0.445 & 0.976 \\
 & LOD/$\sqrt{2}$ & 0.897 & 0.458 & 0.791 & 0.445 & 0.477 & 0.976 \\
 & LOD & 0.896 & 0.449 & 0.797 & 0.465 & 0.446 & 0.967 \\
 & MICE & 0.811 & 0.285 & 0.768 & 0.353 & 0.073 & 0.673 \\
\midrule
\multirow{5}{*}{50\%}
 & Zero & 0.823 & 0.264 & 0.632 & 0.343 & 0.111 & 0.871 \\
 & LOD/2 & 0.896 & 0.385 & 0.781 & 0.433 & 0.370 & 0.951 \\
 & LOD/$\sqrt{2}$ & 0.894 & 0.396 & 0.793 & 0.460 & 0.350 & 0.947 \\
 & LOD & 0.889 & 0.408 & 0.798 & 0.475 & 0.251 & 0.928 \\
 & MICE & 0.289 & 0.065 & 0.654 & 0.241 & 0.005 & 0.515 \\
\bottomrule
\end{tabular}
\end{table}

Key findings:

\begin{itemize}
\item \textbf{Low censoring (5--10\%):} LOD/$\sqrt{2}$, LOD/2, and LOD substitution all yielded comparable results (power $\sim$0.90, accuracy $\sim$0.98). MICE performed comparably at 5\% but showed slight degradation in clustering ARI (0.41 vs.\ 0.51) and accuracy (0.94 vs.\ 0.98) relative to LOD/$\sqrt{2}$. Zero substitution was the clear outlier, with substantially degraded power (0.74 at 5\%, 0.68 at 10\%) and clustering ARI (0.25 and 0.12, respectively).

\item \textbf{Moderate censoring (20--30\%):} Simple substitution methods (LOD/$\sqrt{2}$, LOD/2, LOD) maintained stable performance (power $\sim$0.90, accuracy $\geq$0.97). MICE imputation degraded substantially as censoring increased: at 30\% censoring, accuracy fell to 0.67 and clustering ARI to 0.07, reflecting insufficient observed data for reliable multivariate imputation. Zero substitution remained consistently poor (accuracy 0.87--0.88).

\item \textbf{High censoring (50\%):} MICE imputation collapsed (power 0.29, accuracy 0.52, ARI 0.005), as half the input data consisted of NaN values, providing too little observed information for the iterative regression framework. In contrast, LOD/$\sqrt{2}$ (power 0.89, accuracy 0.95) and LOD/2 (power 0.90, accuracy 0.95) remained robust. Zero substitution (power 0.82, accuracy 0.87) outperformed MICE at this extreme level.

\item \textbf{Normalization method stability:} The relative ranking of normalization methods was remarkably stable across all imputation strategies. CLR consistently maximized power and clustering regardless of LOD handling; its elevated Type~I error persisted across all imputation methods and censoring levels.

\item \textbf{Practical implications:} LOD/$\sqrt{2}$ substitution (Hornung and Reed, 1990) provides robust performance across all censoring levels tested (5--50\%) and is recommended as the default for BAL biomarker studies. Zero substitution should be avoided as it consistently degrades performance. MICE, while theoretically appealing, is not recommended for high-censoring scenarios ($\geq$20\%) in this context, as the proportion of missing values exceeds the capacity of multivariate imputation algorithms. Normalization method choice has a substantially larger effect on analytical conclusions than LOD imputation method choice.
\end{itemize}

% ============================================================
\section{Supplementary Figures}
% ============================================================

\subsection*{Figure~E1: High-Dimensional Analysis}
\textit{Performance metrics as a function of the number of biomarkers ($p \in \{10, 50, 100, 500\}$) under moderate dilution. (A)~Statistical power with and without FDR correction. (B)~Type~I error with and without FDR correction. (C)~Computation time scaling.}

\subsection*{Figure~E2: CLR Type~I Error Decomposition}
\textit{(A)~Type~I error rates under null conditions for each normalization method, comparing standard $t$-tests (red) vs.\ permutation tests (blue). Dashed line indicates nominal $\alpha = 0.05$. (B)~QQ-plots of $p$-values from CLR-transformed data under null conditions: $t$-test $p$-values (left) vs.\ permutation $p$-values (right).}

\subsection*{Figure~E3: Batch--Dilution Interaction}
\textit{Performance metrics under combined batch effects and dilution for five correction strategies. (A)~Classification accuracy. (B)~Clustering ARI. (C)~Statistical power. (D)~Type~I error.}

\subsection*{Figure~E4: PICFLU Cohort Reanalysis}
\textit{Reanalysis of BAL biomarker data from the PICFLU cohort (Boyd et al., 2020). (A)~Number of significant analytes identified under each normalization method. (B)~Heatmap of significance concordance across normalization methods. (C)~Pairwise Jaccard similarity between significant analyte sets.}

\subsection*{Figure~E5: Multi-Group Analysis}
\textit{Normalization performance with 2, 3, and 4 groups. (A)~Power by number of groups. (B)~Clustering ARI by number of groups. (C)~Classification accuracy by number of groups.}

\subsection*{Figure~E6: LOD Sensitivity Heatmaps}
\textit{Heatmaps of six performance metrics (power, Type~I error, correlation recovery, PCA RV coefficient, clustering ARI, classification accuracy) as a function of LOD severity (5--50\% below LOD) and imputation method (zero, LOD/2, LOD/$\sqrt{2}$, LOD, MICE), averaged across all normalization methods. Color intensity indicates metric value.}

\subsection*{Figure~E7: LOD $\times$ Normalization Interaction}
\textit{Power by normalization method and imputation strategy at each LOD severity level (5\%, 10\%, 20\%, 30\%, 50\%). Heatmaps show power averaged across 20 replications under moderate dilution (Beta(5,5)).}

\subsection*{Figure~E8: Radar Plots of Method-Specific Performance}
\textit{Radar plots summarizing six performance metrics for each normalization method under (A)~mild, (B)~moderate, and (C)~severe dilution.}

% ============================================================
\section{Code Availability}
% ============================================================

The complete simulation framework, analysis scripts, and figure-generating code are available at the project repository under the MIT License. The framework is implemented in Python 3.11 and requires NumPy, SciPy, scikit-learn, pandas, Matplotlib, and Seaborn. All random seeds are fixed for full reproducibility.

Key files:
\begin{itemize}
\item \texttt{biomarker\_dilution\_sim.py}: Core simulation engine (data generation, normalization, analysis, evaluation)
\item \texttt{run\_analysis.py}: Complete analysis pipeline generating all figures and tables
\item \texttt{visualization\_module.py}: Publication-quality figure generation
\item \texttt{reanalysis.py}: PICFLU cohort reanalysis script
\end{itemize}

% ============================================================
% REFERENCES
% ============================================================
\bibliographystyle{unsrtnat}
\bibliography{references}

\end{document}
